\documentclass[
  11pt,
  a4paper,
  oneside,
  parskip=half,
  bibliography=totoc,
  listof=totoc,
  numbers=noenddot
]{scrreprt}

% ---------------------------------------------------------------
% Angaben zur Arbeit
% ---------------------------------------------------------------
\newcommand{\arbeitstitel}{Medcheck}
\newcommand{\arbeitsuntertitel}{Entwicklung und Evaluation eines LLM-gestützten Arzneimittel-Interaktionsprüfers auf Basis amtlicher FDA-Fachinformationen}
\newcommand{\autor}{Vorname Nachname}
\newcommand{\matrikel}{Matrikelnummer}
\newcommand{\hochschule}{Name der Hochschule}
\newcommand{\studiengang}{Bachelorstudiengang Informatik}
\newcommand{\erstbetreuer}{Erstbetreuer:in}
\newcommand{\zweitbetreuer}{Zweitbetreuer:in}
\newcommand{\abgabedatum}{Oktober 2026}

% ---------------------------------------------------------------
% Pakete
% ---------------------------------------------------------------
\usepackage[T1]{fontenc}
\usepackage[utf8]{inputenc}
\usepackage[ngerman]{babel}
\usepackage{lmodern}
\usepackage{microtype}
\usepackage{csquotes}

% Satzspiegel und Zeilenabstand: bei abweichenden Vorgaben der Hochschule
% hier anpassen (z. B. \setstretch{1.5}); der Seitenumfang ändert sich dann.
\usepackage[left=3cm, right=2.5cm, top=2.5cm, bottom=2.5cm]{geometry}
\usepackage{setspace}
\setstretch{1.0}

% Kompaktere Überschriften
\RedeclareSectionCommand[beforeskip=0.5\baselineskip, afterskip=1.2\baselineskip, tocbeforeskip=0.5em]{chapter}
\RedeclareSectionCommand[beforeskip=-2.6ex plus -1ex minus -.2ex, afterskip=1.2ex plus .2ex]{section}
\RedeclareSectionCommand[beforeskip=1.2ex plus .6ex minus .2ex]{paragraph}
\setkomafont{chapter}{\usekomafont{disposition}\LARGE}

\usepackage[
  backend=biber,
  style=numeric,
  sorting=none,
  maxbibnames=6,
  giveninits=true,
  urldate=long
]{biblatex}
\addbibresource{literatur.bib}
\renewcommand*{\bibfont}{\small}
\setcounter{biburlnumpenalty}{9000}
\setcounter{biburllcpenalty}{7000}
\setcounter{biburlucpenalty}{8000}
\biburlbigskip=0mu

% Etwas mehr Spielraum für Zeilen mit langen Code-Bezeichnern
\setlength{\emergencystretch}{1.5em}

\usepackage{graphicx}
\usepackage{booktabs}
\usepackage{tabularx}
\usepackage{array}
\newcolumntype{L}{>{\raggedright\arraybackslash}X}
\usepackage{multirow}
\usepackage[table]{xcolor}
\usepackage{enumitem}
\setlist{itemsep=0.15em, topsep=0.3em}

\usepackage{tikz}
\usetikzlibrary{arrows.meta, positioning, fit, backgrounds}
\usepackage{pgfplots}
\pgfplotsset{compat=1.18}

\usepackage{caption}
\captionsetup{font=small, labelfont=bf, format=plain}

\usepackage{listings}
\definecolor{codebg}{HTML}{F3F6F7}
\definecolor{codeframe}{HTML}{C9D3D7}
\definecolor{codekw}{HTML}{1C5A70}
\definecolor{codecomment}{HTML}{6B7A82}
\definecolor{codestring}{HTML}{8A4B08}
\lstset{
  basicstyle=\ttfamily\footnotesize,
  backgroundcolor=\color{codebg},
  frame=single,
  rulecolor=\color{codeframe},
  framerule=0.4pt,
  keywordstyle=\color{codekw}\bfseries,
  commentstyle=\color{codecomment}\itshape,
  stringstyle=\color{codestring},
  breaklines=true,
  columns=fullflexible,
  keepspaces=true,
  showstringspaces=false,
  aboveskip=0.8em,
  belowskip=0.4em,
  xleftmargin=0.4em,
  xrightmargin=0.4em,
  literate={ä}{{\"a}}1 {ö}{{\"o}}1 {ü}{{\"u}}1 {ß}{{\ss}}1
           {Ä}{{\"A}}1 {Ö}{{\"O}}1 {Ü}{{\"U}}1 {–}{{--}}1
}
\renewcommand{\lstlistingname}{Listing}
\renewcommand{\lstlistlistingname}{Listingverzeichnis}

\usepackage[hidelinks]{hyperref}
\usepackage{bookmark}

% Farben der Risikostufen (identisch zur Benutzeroberfläche)
\definecolor{riskhigh}{HTML}{C62828}
\definecolor{riskmid}{HTML}{EF6C00}
\definecolor{risklow}{HTML}{2E7D32}
\definecolor{varA}{HTML}{8A979E}
\definecolor{varB}{HTML}{1C5A70}
\definecolor{orderA}{HTML}{5B6B73}
\definecolor{orderB}{HTML}{D08C2E}
\newcommand{\riskdot}[1]{\tikz[baseline=-0.6ex]\fill[#1] (0,0) circle (0.8ex);}
\newcommand{\riskoff}{\tikz[baseline=-0.6ex]\draw[gray, line width=0.5pt] (0,0) circle (0.75ex);}

% Makros für den Anhang mit den Rohdaten der Befragung
\newenvironment{answers}{\begin{itemize}[leftmargin=1.2em, label={}, itemsep=0.1em]}{\end{itemize}}
\newcommand{\qa}[2]{\item #1 \hfill \textbf{#2}}
\newcommand{\block}[1]{\item[] \textsc{#1}}

\hypersetup{
  pdftitle={\arbeitstitel: \arbeitsuntertitel},
  pdfauthor={\autor}
}

\begin{document}

% ---------------------------------------------------------------
% Titelseite
% ---------------------------------------------------------------
\begin{titlepage}
  \centering
  {\large \hochschule\par}
  \vspace{0.4em}
  {\normalsize \studiengang\par}
  \vspace{4.5cm}
  {\sffamily\bfseries\Huge \arbeitstitel\par}
  \vspace{1em}
  {\sffamily\Large \arbeitsuntertitel\par}
  \vspace{2.5cm}
  {\large Bachelorarbeit\par}
  \vspace{0.6em}
  {\normalsize zur Erlangung des akademischen Grades\\ Bachelor of Science (B.\,Sc.)\par}
  \vfill
  \begin{tabular}{@{}ll@{}}
    Vorgelegt von: & \autor \\
    Matrikelnummer: & \matrikel \\
    Erstbetreuung: & \erstbetreuer \\
    Zweitbetreuung: & \zweitbetreuer \\
    Abgabe: & \abgabedatum
  \end{tabular}
\end{titlepage}

\pagenumbering{roman}
\addchap*{Abstract}

Wechselwirkungen zwischen Arzneimitteln sind eine häufige und vermeidbare Ursache unerwünschter Arzneimittelereignisse. Die maßgeblichen Informationen sind zwar öffentlich zugänglich, liegen jedoch überwiegend als fachsprachliche, englische Fachinformationen vor, die für medizinische Laien kaum verständlich sind. Gleichzeitig zeigen aktuelle Studien, dass Large Language Models (LLMs) Wechselwirkungen zwar mit hoher Sensitivität, aber geringer Spezifität erkennen, wenn sie aus ihrem eigenen Trainingswissen antworten.

Diese Arbeit stellt mit \emph{Medcheck} einen Prototyp vor, der beide Beobachtungen verbindet: Das Sprachmodell wird nicht als Wissensquelle, sondern ausschließlich als Erklärwerkzeug über amtliche Daten eingesetzt. Frei eingegebene Medikamentennamen werden über RxNorm normalisiert, der Abschnitt \enquote{Drug Interactions} der Fachinformation wird aus der openFDA Drug Label API bezogen, und ein austauschbares LLM (Anthropic Claude, Google Gemini oder ein lokales Modell über Ollama) überführt diesen Text in eine dreistufige Risikoeinstufung und eine laienverständliche Erklärung in einer von elf Sprachen. Die Implementierung auf Basis von Spring Boot und Angular legt besonderen Wert auf Anbieterunabhängigkeit und auf kontrollierte Degradation bei Ausfall externer Dienste.

In einer Befragung mit elf Teilnehmenden wurde die von Medcheck erzeugte Erklärung zur Kombination von Warfarin und Ibuprofen mit dem entsprechenden Text von Drugs.com verglichen. Die Medcheck-Variante erzielte im Mittel höhere Werte für Verständlichkeit (4{,}30 gegenüber 3{,}64 auf einer fünfstufigen Skala) und wurde von acht der elf Personen als verständlicher bewertet. Beim Vertrauen lagen beide Varianten nahezu gleichauf; mehrere Teilnehmende vermissten in der kurzen Erklärung konkrete Warnzeichen. Die Unterschiede sind bei dieser Stichprobengröße statistisch nicht signifikant. Die Ergebnisse sprechen für eine kombinierte Darstellung aus knapper Zusammenfassung und aufklappbaren Detailinformationen, wie sie der Prototyp bereits vorsieht.


\tableofcontents

\clearpage
\pagenumbering{arabic}

\chapter{Einleitung}
\label{ch:einleitung}

\section{Motivation}

Mit steigendem Alter und zunehmender Zahl chronischer Erkrankungen nehmen viele Menschen mehrere Arzneimittel gleichzeitig ein. Diese Polypharmazie erhöht das Risiko von Arzneimittelwechselwirkungen (engl.\ \emph{drug--drug interactions}, DDIs), bei denen ein Wirkstoff die Wirkung oder die Verstoffwechselung eines anderen verändert. Ein klassisches Beispiel ist die Kombination des Gerinnungshemmers Warfarin mit dem nichtsteroidalen Antirheumatikum Ibuprofen: Beide Wirkstoffe beeinträchtigen die Blutgerinnung, sodass das Risiko schwerer, insbesondere gastrointestinaler Blutungen deutlich steigt. Ibuprofen ist zudem rezeptfrei erhältlich und wird daher häufig ohne ärztliche Rücksprache eingenommen.

Für Fachpersonal existieren etablierte Interaktionsdatenbanken und klinische Entscheidungsunterstützungssysteme (CDSS), etwa Lexicomp oder UpToDate~\cite{UpToDate}. Für Laien stehen frei zugängliche Interaktionschecker wie Drugs.com~\cite{DrugsCom2025} oder der Wechselwirkungs-Check der Apotheken Umschau~\cite{ApothekenUmschau} zur Verfügung. Studien zeigen jedoch, dass sich die Ergebnisse verschiedener Datenbanken unterscheiden~\cite{article, Roca2022} und dass die Texte für medizinische Laien oft schwer verständlich sind. Gleichzeitig nutzen immer mehr Menschen KI-Assistenten wie Gemini~\cite{GoogleGemini} für Gesundheitsfragen. Deren Antworten sind gut lesbar, beruhen aber auf dem Trainingswissen des Modells und sind nicht ohne Weiteres überprüfbar. Vergleichsstudien berichten für solche Modelle eine hohe Sensitivität bei gleichzeitig geringer Spezifität, also eine Tendenz, Wechselwirkungen zu überschätzen~\cite{MOHAMMED2026100733}.

\section{Problemstellung und Zielsetzung}

Es besteht damit eine Lücke zwischen verlässlichen, aber schwer verständlichen Fachinformationen einerseits und gut verständlichen, aber nicht ausreichend verlässlichen KI-Antworten andererseits. Ziel dieser Arbeit ist die Entwicklung und Evaluation eines Prototyps, der diese Lücke schließt. Der Prototyp \emph{Medcheck} bezieht die Interaktionsinformationen aus der amtlichen Datenquelle openFDA~\cite{openFDA} und setzt ein Sprachmodell ausschließlich dazu ein, diese Informationen in eine verständliche Form zu bringen. Das Modell soll dabei weder eigenes medizinisches Wissen hinzufügen noch individuelle Therapieempfehlungen geben.

Daraus ergeben sich zwei Forschungsfragen:

\begin{description}[leftmargin=2.6em, style=sameline, font=\normalfont\bfseries]
  \item[FF1] Wie lässt sich ein Interaktionsprüfer technisch realisieren, der amtliche Fachinformationen mit einem austauschbaren Sprachmodell zu laienverständlichen Erklärungen verarbeitet und dabei robust gegenüber dem Ausfall externer Dienste bleibt?
  \item[FF2] Wie bewerten medizinische Laien die Verständlichkeit, das Vertrauen und die Nützlichkeit einer so erzeugten Erklärung im Vergleich zu einem etablierten Interaktionschecker?
\end{description}

\section{Aufbau der Arbeit}

Kapitel~\ref{ch:sota} stellt bestehende Interaktionsdatenbanken, offene Datenquellen und den Forschungsstand zu LLMs bei Wechselwirkungen dar. Kapitel~\ref{ch:methodik} beschreibt Vorgehen, Anforderungen und Befragungsdesign. Kapitel~\ref{ch:implementierung} stellt die Implementierung vor und beantwortet FF1, Kapitel~\ref{ch:evaluation} präsentiert und diskutiert die Befragung zu FF2. Kapitel~\ref{ch:konklusion} fasst die Ergebnisse zusammen und gibt einen Ausblick.

\chapter{State of the Art}
\label{ch:sota}

\section{Interaktionsdatenbanken und klinische Entscheidungsunterstützung}

Die Prüfung auf Arzneimittelwechselwirkungen stützt sich in der Praxis auf kuratierte Datenbanken. Für Fachpersonal haben sich kommerzielle Systeme wie Lexicomp und UpToDate etabliert, die Wechselwirkungen nach Schweregrad klassifizieren und mit Angaben zu Mechanismus und Management versehen~\cite{UpToDate}. Daneben existieren frei zugängliche Angebote, die sich an Fachpersonal und Laien richten. Tabelle~\ref{tab:werkzeuge} ordnet die in dieser Arbeit betrachteten Werkzeuge ein.

\begin{table}[htbp]
  \centering
  \small
  \caption{Auswahl bestehender Werkzeuge und Informationsquellen zu Arzneimittelwechselwirkungen.}
  \label{tab:werkzeuge}
  \begin{tabularx}{\textwidth}{@{}l L L@{}}
    \toprule
    Werkzeug & Art & Primäre Zielgruppe \\
    \midrule
    UpToDate / Lexicomp~\cite{UpToDate} & Kommerzielles CDSS, Point-of-Care & Ärztliches und pharmazeutisches Fachpersonal \\
    Medscape~\cite{Medscape} & Medizinisches Fachportal mit Interaktionschecker & Fachpersonal \\
    DrugBank~\cite{DrugBank} & Pharmakologische Wissensdatenbank & Forschung, Industrie, Fachpersonal \\
    Drugs.com~\cite{DrugsCom2025} & Frei zugänglicher Interaktionschecker & Laien und Fachpersonal \\
    WebMD~\cite{WebMD} & Gesundheitsportal & Laien \\
    Patient.info~\cite{PatientInfo} & Gesundheitsportal & Laien \\
    Apotheken Umschau~\cite{ApothekenUmschau} & Deutschsprachiger Wechselwirkungs-Check & Laien \\
    Drug Interaction Tool~\cite{DrugInteractionTool} & Webbasierter Interaktionschecker & Laien \\
    Gemini~\cite{GoogleGemini} & Allgemeiner KI-Assistent & Allgemeine Öffentlichkeit \\
    \bottomrule
  \end{tabularx}
\end{table}

Ein grundlegendes Problem ist die mangelnde Übereinstimmung zwischen den Quellen. Nabeel und Birand untersuchten potenzielle Wechselwirkungen in öffentlichen Apotheken in Nordzypern mit drei verschiedenen Datenbanken und verglichen deren Ergebnisse~\cite{article}. Roca und Roca bewerteten frei verfügbare elektronische Interaktionschecker anhand multimorbider Patientinnen und Patienten~\cite{Roca2022}. Beide Arbeiten verdeutlichen, dass Anzahl und Einstufung erkannter Wechselwirkungen von der gewählten Quelle abhängen. Für die Nutzerinnen und Nutzer bedeutet das, dass das Ergebnis einer einzelnen Prüfung nicht als absolute Wahrheit verstanden werden sollte.

In klinischen Informationssystemen tritt ein weiteres Problem auf. Interaktionswarnungen werden sehr häufig übersteuert, weil viele Meldungen als klinisch nicht relevant wahrgenommen werden. Eine Scoping-Review von Villa Zapata et al.\ fasst die Literatur zu diesem Phänomen zusammen~\cite{VillaZapata}, das unter dem Begriff \emph{Alert Fatigue} bekannt ist. Van De Sijpe et al.\ evaluierten die Gesamtleistung eines DDI-CDSS quantitativ und ergänzten dies um eine Befragung der Anwenderinnen und Anwender~\cite{VanDeSijpe}. Für diese Arbeit ist daraus abzuleiten, dass ein Werkzeug nicht nur korrekt, sondern auch spezifisch sein muss: Zu viele oder zu drastische Warnungen untergraben die Akzeptanz.

\section{Offene Datenquellen}

Zwei frei zugängliche Datenquellen der US-Bundesbehörden bilden die Datengrundlage des in dieser Arbeit entwickelten Prototyps. Die \emph{openFDA}-Initiative der U.S.\ Food and Drug Administration stellt seit 2014 öffentliche Datensätze der Behörde über REST-Schnittstellen bereit, darunter die Fachinformationen aller in den USA zugelassenen Arzneimittel im Format \emph{Structured Product Labeling} (SPL)~\cite{openFDA}. Abschnitt~7 dieser Fachinformationen beschreibt die Wechselwirkungen. \emph{RxNav} ist das Werkzeugangebot der National Library of Medicine rund um das Vokabular RxNorm, das klinische Arzneimittelbezeichnungen normalisiert und jedem Konzept einen eindeutigen Bezeichner (RxCUI) zuordnet~\cite{RxNav}. RxNav bietet unter anderem eine unscharfe Suche, die auch fehlerhaft geschriebene Eingaben einem Wirkstoff zuordnet. Beide Quellen sind ohne Lizenzkosten nutzbar, liefern aber keine laienverständlichen Erklärungen.

\section{Maschinelles Lernen zur Vorhersage von Wechselwirkungen}

Vor dem Aufkommen großer Sprachmodelle wurde maschinelles Lernen vor allem zur \emph{Vorhersage} unbekannter Wechselwirkungen eingesetzt. Chen et al.\ geben einen systematischen Überblick über KI-Verfahren zur Vorhersage von Arzneimittel-Arzneimittel-, Arzneimittel-Nahrungsmittel- und Arzneimittel-Mikrobiom-Interaktionen~\cite{10.1093/bib/bbac427}. Unter den Modellen für Arzneimittelwechselwirkungen heben die Autorinnen und Autoren Ansätze hervor, die auf dem Metabolismus über Cytochrom-P450-Enzyme, auf der Ähnlichkeit von Wirkstoffen oder auf deren Zielstrukturen beruhen. Diese Verfahren zielen auf Arzneimittelentwicklung und Sicherheitsüberwachung, nicht auf die Kommunikation bekannter Risiken an Patientinnen und Patienten. Sie sind daher für die vorliegende Fragestellung nur mittelbar relevant.

\section{Large Language Models zur Erkennung von Wechselwirkungen}

Mit der Verbreitung von LLMs rückte die Frage in den Vordergrund, ob allgemeine Sprachmodelle bekannte Wechselwirkungen zuverlässig erkennen. Bischof et al.\ verglichen ChatGPT mit etablierten CDSS bei der Analyse von Arzneimittelwechselwirkungen~\cite{Bischof2025}. Sicard et al.\ untersuchten, ob LLMs Wechselwirkungen erkennen, die zu unerwünschten Arzneimittelwirkungen geführt hatten~\cite{Sicard2025}. Blotske et al.\ sowie Tilley et al.\ evaluierten LLMs für die Identifikation von Wechselwirkungen und betonen dabei die Bedeutung des Umgangs mit Unsicherheit in den Modellantworten~\cite{Blotske2025, Tilley2026}.

Besonders aufschlussreich ist die Querschnittsstudie von Mohammed et al.~\cite{MOHAMMED2026100733}. Die Autorinnen und Autoren prüften 186 Kombinationen von Antiepileptika mit weiteren Arzneimitteln, von denen 126 nach Lexicomp als schwer oder mittelschwer eingestuft sind. Lexicomp diente als Referenzstandard, verglichen wurden ChatGPT-4, DeepSeek-V3 und Drugs.com. Tabelle~\ref{tab:mohammed} fasst die Ergebnisse zusammen.

\begin{table}[htbp]
  \centering
  \small
  \caption{Erkennung klinisch relevanter Wechselwirkungen im Vergleich zu Lexicomp nach Mohammed et al.~\cite{MOHAMMED2026100733}.}
  \label{tab:mohammed}
  \begin{tabular}{@{}lccc@{}}
    \toprule
    System & Sensitivität & Spezifität & Spezifität nach iterativem Prompting \\
    \midrule
    Drugs.com   & 0{,}870 & 0{,}629 & -- \\
    ChatGPT-4   & 0{,}842 & 0{,}358 & 0{,}77 \\
    DeepSeek-V3 & 0{,}877 & 0{,}200 & 0{,}42 \\
    \bottomrule
  \end{tabular}
\end{table}

Beide Sprachmodelle erkannten nahezu so viele relevante Wechselwirkungen wie Drugs.com, stuften aber deutlich häufiger harmlose oder geringfügige Kombinationen als klinisch relevant ein. Die Autorinnen und Autoren weisen darauf hin, dass diese Überschätzung zu Alert Fatigue sowie zu unnötiger Überwachung oder Therapieänderung führen kann. Ein strukturiertes, an Evidenz gebundenes Nachfragen (\emph{iterative prompting}) verbesserte die Spezifität bei den falsch-positiven Fällen erheblich. Drugs.com zeigte insgesamt die ausgewogenste Leistung (Genauigkeit 0{,}709). Die Studie schließt, dass LLMs derzeit als ergänzende Erklärwerkzeuge dienen können, strukturierte Interaktionsdatenbanken für das eigentliche Screening aber verlässlicher bleiben.

Diese Einschätzung deckt sich mit übergreifenden Arbeiten. Alsayed analysiert Zuverlässigkeit und Risiken KI-gestützter Systeme für Medikationsentscheidungen und die Folgen fehlerhafter Ausgaben~\cite{alsayed2026aigetswrongreliability}. Die Scoping-Review von Ong et al.\ gibt einen Überblick über den Einsatz generativer KI und großer Sprachmodelle zur Vermeidung medikationsbezogener Schäden~\cite{Ong2025}.

\section{Sprachmodelle als Erklärwerkzeug für Patientinnen und Patienten}

Neben der Erkennung von Wechselwirkungen wird untersucht, ob generative KI Arzneimittelinformationen für Patientinnen und Patienten verständlicher machen kann. Almeida et al.\ berichten erste Ergebnisse zu personalisierten, mit generativer KI erstellten Packungsbeilagen~\cite{Almeida2025}. Ansätze dieser Art nutzen die sprachlichen Stärken der Modelle, ohne ihnen die Rolle einer medizinischen Wissensquelle zuzuweisen.

\section{Forschungslücke}

Aus dem Forschungsstand ergeben sich drei Beobachtungen. Erstens erkennen LLMs Wechselwirkungen aus ihrem Trainingswissen mit unzureichender Spezifität~\cite{MOHAMMED2026100733}. Zweitens sind verlässliche Datenquellen vorhanden, aber nicht laienverständlich aufbereitet~\cite{openFDA}. Drittens liegt der Schwerpunkt bisheriger Evaluationen auf diagnostischen Kennzahlen gegenüber einem Referenzstandard; wie Laien die Erklärungen wahrnehmen, ist deutlich seltener untersucht. Diese Arbeit setzt an diesen Punkten an. Sie bindet das Sprachmodell konsequent an den amtlichen Text (\emph{Grounding}), anstatt es frei antworten zu lassen, und evaluiert das Ergebnis aus der Perspektive medizinischer Laien.

\chapter{Methodik}
\label{ch:methodik}

\section{Vorgehensmodell}

Die Arbeit folgt einem gestaltungsorientierten Vorgehen in zwei Phasen. In der ersten Phase wurde ein Prototyp iterativ entwickelt: Jede Iteration umfasste die Umsetzung einer Teilfunktion, den Test gegen die realen Schnittstellen von RxNorm, openFDA und der LLM-Anbieter sowie die Korrektur der aufgetretenen Abweichungen. Dies war notwendig, weil sich das Verhalten der externen Schnittstellen in mehreren Punkten von ihrer Dokumentation unterschied (Abschnitt~\ref{sec:herausforderungen}). In der zweiten Phase wurde eine vom Prototyp erzeugte Erklärung in einer Befragung mit medizinischen Laien mit einem etablierten Interaktionschecker verglichen.

\section{Anforderungen}

Aus Problemstellung und Forschungsstand wurden die Anforderungen in Tabelle~\ref{tab:anforderungen} abgeleitet. Die nichtfunktionalen Anforderungen NF1 bis NF3 adressieren unmittelbar die in Kapitel~\ref{ch:sota} beschriebenen Schwächen: die Abhängigkeit von Drittanbietern, die schnelle Weiterentwicklung des LLM-Marktes und die geringe Spezifität frei antwortender Modelle.

\begin{table}[htbp]
  \centering
  \small
  \caption{Funktionale (F) und nichtfunktionale (NF) Anforderungen an den Prototyp.}
  \label{tab:anforderungen}
  \begin{tabularx}{\textwidth}{@{}l L L@{}}
    \toprule
    ID & Anforderung & Umsetzung \\
    \midrule
    F1  & Tippfehlertolerante Suche nach Medikamenten & RxNorm-Suche, Autovervollständigung \\
    F2  & Prüfung von zwei bis zwanzig Medikamenten & Dynamische Eingabefelder, Validierung \\
    F3  & Dreistufige Risikoeinstufung & LLM-Klassifikation, Ampel \\
    F4  & Laienverständliche Erklärung, wählbare Sprache & Parametrisierte Prompts, elf Sprachen \\
    F5  & Einsicht in den unveränderten Originaltext & Umschaltbare Ansicht \enquote{FDA-Original} \\
    NF1 & Robustheit bei Ausfall externer Dienste & Kontrollierte Degradation, Timeouts \\
    NF2 & Austauschbarkeit des LLM-Anbieters & Interface \texttt{LlmProvider}, Konfiguration \\
    NF3 & Nachvollziehbarkeit der Aussagen & Grounding auf openFDA-Text \\
    NF4 & Geringe Betriebskosten & Caching, gekürzte Eingabe, lokales Modell \\
    \bottomrule
  \end{tabularx}
\end{table}

\section{Entwurfsprinzipien}

Drei Prinzipien leiten den Systementwurf. \emph{Grounding}: Das Sprachmodell erhält den amtlichen Text als einzige Wissensgrundlage und wird angewiesen, nichts hinzuzufügen. Damit wird die von Mohammed et al.\ beschriebene Überschätzung von Wechselwirkungen durch frei antwortende Modelle~\cite{MOHAMMED2026100733} strukturell begrenzt. \emph{Anbieterunabhängigkeit}: Die LLM-Schicht ist hinter einer Schnittstelle gekapselt, sodass Modelle und Anbieter ohne Codeänderung austauschbar sind. \emph{Kontrollierte Degradation}: Fällt ein externer Dienst aus, liefert das System die verfügbaren Informationen aus, statt abzubrechen; insbesondere bleibt der amtliche Text sichtbar, wenn das Sprachmodell nicht antwortet.

\section{Design der Nutzerbefragung}
\label{sec:befragungsdesign}

\paragraph{Untersuchungsgegenstand.} Verglichen wurden zwei Erklärungen zur Kombination von Warfarin und Ibuprofen. \emph{Variante~A} ist der Erklärungstext des Interaktionscheckers von Drugs.com~\cite{DrugsCom2025}. Er beschreibt das Blutungsrisiko, empfiehlt Rücksprache mit dem Arzt, erwähnt mögliche Alternativen, Dosisanpassung und INR-Kontrolle und zählt konkrete Warnzeichen einer Blutung auf, etwa Schwindel, teerartigen Stuhl oder kaffeesatzartiges Erbrechen. Für Teilnehmende mit Schwierigkeiten im Englischen lag eine maschinell (mit Claude) erstellte deutsche Übersetzung vor. \emph{Variante~B} ist die von Medcheck erzeugte Zusammenfassung in zwei bis drei Sätzen. Sie benennt das erhöhte Blutungsrisiko und den synergistischen Mechanismus und fordert zu ärztlicher Rücksprache und Überwachung auf. Ihre englische und deutsche Fassung wurden über die Spracheinstellung des Prototyps erzeugt.

\paragraph{Ablauf.} Die Befragung folgte einem Within-Subject-Design: Jede Person las und bewertete beide Varianten, wobei die Reihenfolge zwischen den Personen variiert wurde. Teil~0 erfasste Altersgruppe, Einverständnis, regelmäßige Medikamenteneinnahme, bisherige Online-Recherche und die Sicherheit im Umgang mit medizinischen Informationen. In Teil~1 und~2 wurden die Varianten A und B anhand derselben acht Aussagen bewertet. Teil~3 verglich beide Varianten direkt in vier Fragen (verständlicher, vertrauenswürdiger, hilfreicher für eine Entscheidung, eher weiterzuempfehlen) mit offener Begründung. Teil~4 war ein Verständnistest mit zwei offenen Fragen: was als Nächstes zu tun sei und wie gefährlich die Kombination sei. Teil~5 enthielt offene Fragen zu Gefallen, Störendem, fehlenden Informationen und möglichen Nutzungssituationen.

\paragraph{Messinstrument.} Die acht Aussagen aus Teil~1 und~2 wurden auf einer fünfstufigen Likert-Skala von 1 (\enquote{trifft gar nicht zu}) bis 5 (\enquote{trifft voll zu}) bewertet und drei Dimensionen zugeordnet (Tabelle~\ref{tab:items}). Die negativ formulierte Aussage V3 wurde für Dimensions- und Gesamtwerte umkodiert (\(x' = 6 - x\)).

\begin{table}[htbp]
  \centering
  \small
  \caption{Bewertungsaussagen der Teile~1 und~2.}
  \label{tab:items}
  \begin{tabularx}{\textwidth}{@{}l l L@{}}
    \toprule
    Dimension & Kürzel & Aussage \\
    \midrule
    \multirow{3}{*}{Verständlichkeit} & V1 & Die Erklärung war für mich leicht zu verstehen. \\
                                      & V2 & Die verwendeten Begriffe waren mir klar. \\
                                      & V3 & Ich musste Teile mehrfach lesen, um sie zu verstehen. (negativ) \\
    \midrule
    \multirow{2}{*}{Vertrauen}        & T1 & Ich vertraue der Richtigkeit dieser Information. \\
                                      & T2 & Die Erklärung wirkte seriös und glaubwürdig. \\
    \midrule
    \multirow{3}{*}{Nützlichkeit}     & N1 & Die Erklärung half mir, das Risiko einzuschätzen. \\
                                      & N2 & Nach dem Lesen weiß ich, was ich als Nächstes tun sollte. \\
                                      & N3 & Die Erklärung enthielt alle Informationen, die ich brauchte. \\
    \bottomrule
  \end{tabularx}
\end{table}

\paragraph{Durchführung und Stichprobe.} Die Befragung fand zwischen dem 18.~August und dem 10.~September 2026 statt. Teilgenommen haben elf Personen, die über das persönliche Umfeld gewonnen wurden (Gelegenheitsstichprobe). Alle Personen willigten in die Aufzeichnung der Daten für diese Arbeit ein. Ein Teil der Befragungen wurde mündlich geführt und transkribiert.\IfFileExists{kapitel/anhang-umfrage.tex}{ Die Rohdaten und Transkripte sind im Anhang dokumentiert.}{}

\paragraph{Auswertung.} Je Aussage und Variante werden Mittelwert, Standardabweichung und Median berichtet. Dimensions- und Gesamtwerte werden paarweise mit dem exakten Wilcoxon-Vorzeichen-Rang-Test verglichen, der für kleine Stichproben und ordinalskalierte Daten geeignet ist; Personen mit identischem Wert in beiden Varianten gehen nicht in den Test ein. Als Effektstärke wird die Rang-biseriale Korrelation \(r = (W^{+} - W^{-})/(W^{+} + W^{-})\) angegeben, wobei \(W^{+}\) und \(W^{-}\) die Rangsummen der Differenzen zugunsten von B bzw.\ A bezeichnen. Die Präferenzen aus dem direkten Vergleich werden mit einem zweiseitigen Binomialtest gegen die Gleichverteilung geprüft. Das Signifikanzniveau beträgt \(\alpha = 0{,}05\). Die offenen Antworten und Transkripte werden inhaltlich zu Themen zusammengefasst.

\chapter{Implementierung}
\label{ch:implementierung}

\section{Systemarchitektur}

Medcheck folgt einer dreischichtigen Client-Server-Architektur (Abbildung~\ref{fig:architektur}). Eine Single-Page-Application auf Basis von Angular~17 kommuniziert über eine REST-Schnittstelle mit einem Backend auf Basis von Spring Boot (Java~21). Nur das Backend spricht mit externen Diensten. Dadurch bleiben API-Schlüssel außerhalb des Browsers, Ergebnisse können serverseitig zwischengespeichert werden, und Fehlerbehandlung sowie Protokollierung erfolgen an einer zentralen Stelle. Als Zwischenspeicher dient eine In-Memory-Datenbank H2.

\begin{figure}[htbp]
  \centering
  \begin{tikzpicture}[
      node distance=0.55cm and 1.1cm,
      box/.style={draw=varB, line width=0.6pt, rounded corners=2pt, align=center,
                  font=\small\sffamily, minimum height=0.95cm, inner sep=4pt, fill=white},
      ext/.style={box, draw=varA, fill=codebg},
      arr/.style={-{Stealth[length=2.2mm]}, line width=0.6pt, draw=black!70},
      lbl/.style={font=\scriptsize\sffamily, fill=white, inner sep=1.5pt}
    ]
    \node[box, minimum width=2.6cm] (ui) {Browser\\Angular-SPA};
    \node[box, right=1.9cm of ui, minimum width=3.0cm, minimum height=2.6cm] (be) {Spring-Boot-\\Backend};
    \node[box, below=0.7cm of be, minimum width=1.8cm] (h2) {H2-Cache};
    \node[ext, right=1.6cm of be, yshift=1.05cm, minimum width=2.9cm] (rx) {RxNorm API (NLM)};
    \node[ext, right=1.6cm of be, minimum width=2.9cm] (fda) {openFDA Drug Label API};
    \node[ext, right=1.6cm of be, yshift=-1.05cm, minimum width=2.9cm] (llm) {LLM-Anbieter\\{\scriptsize Claude $\mid$ Gemini $\mid$ Ollama}};
    \draw[arr] (ui) -- node[lbl, above=1pt] {REST/JSON} (be);
    \draw[arr] (be) -- (h2);
    \draw[arr] (be.east |- rx) -- node[lbl, above=1pt] {Suche} (rx);
    \draw[arr] (be.east |- fda) -- node[lbl, above=1pt] {Fachinfo} (fda);
    \draw[arr] (be.east |- llm) -- node[lbl, above=1pt] {Prompt} (llm);
  \end{tikzpicture}
  \caption{Komponentenübersicht von Medcheck. Zur Laufzeit ist genau einer der drei LLM-Anbieter aktiv.}
  \label{fig:architektur}
\end{figure}

Das Backend ist nach Verantwortlichkeiten in Pakete gegliedert: \texttt{controller}, \texttt{service}, \texttt{repository}, \texttt{model}, \texttt{dto}, \texttt{validation}, \texttt{exception} und \texttt{config}. JPA-Entitäten verlassen die Service-Schicht nie; an der REST-Grenze werden ausschließlich Data Transfer Objects ausgegeben, die als unveränderliche Java Records umgesetzt sind. Alle Abhängigkeiten werden per Konstruktor injiziert, sodass jede Klasse ohne Spring-Kontext getestet werden kann. Für den Zugriff auf externe Dienste dient der \texttt{WebClient} aus Spring WebFlux, der je Dienst eigene Timeouts und Puffergrößen erlaubt.

Der zentrale Anwendungsfall ist der Aufruf \texttt{POST /api/interactions/check} mit einer Liste von Medikamentennamen und einem Sprachcode. Nach der Validierung löst der \texttt{InteractionService} jeden Namen über RxNorm auf, sucht den Interaktionstext im Cache und fragt bei einem Fehltreffer openFDA ab. Anschließend startet er die Aufrufe an das Sprachmodell parallel: einen zur Bewertung der Kombination und je Medikament einen zur Zusammenfassung der wichtigsten Warnhinweise.

\section{Datenbeschaffung}

\paragraph{Medikamentensuche.} Der \texttt{RxNormService} kapselt den Endpunkt \texttt{approximateTerm} mit \texttt{maxEntries=10}. Die Antwort wird zunächst in Records abgebildet, die das externe JSON-Format spiegeln, und anschließend in das eigene DTO \texttt{RxNormCandidate(rxcui, name, score, rank)} überführt. Kandidaten ohne RxCUI oder Namen werden verworfen, gleichnamige zusammengeführt. Die Trennung zwischen externem Format und eigenem DTO hält die Schnittstelle zum Frontend stabil, auch wenn die NLM ihr Antwortformat ändert.

\paragraph{Abfragekaskade.} Naheliegend wäre, die über RxNorm ermittelte RxCUI direkt in openFDA zu suchen. In der Praxis ist das Feld \texttt{openfda.rxcui} jedoch nur bei einem Teil der Fachinformationen befüllt; für Ibuprofen und Warfarin liefert diese Suche keinen Treffer. Der \texttt{OpenFDAService} probiert deshalb mehrere Suchausdrücke in absteigender Spezifität, bis einer einen Text liefert (Tabelle~\ref{tab:kaskade}). Bei einwortigen Namen wird auf Anführungszeichen verzichtet, weil die exakte Phrasensuche häufig an abweichender Groß- und Kleinschreibung scheitert.

\begin{table}[htbp]
  \centering
  \small
  \caption{Abfragekaskade gegen openFDA am Beispiel Ibuprofen.}
  \label{tab:kaskade}
  \begin{tabularx}{\textwidth}{@{}c l L@{}}
    \toprule
    Stufe & Suchausdruck & Begründung \\
    \midrule
    1 & \texttt{openfda.rxcui:5640} & Eindeutiger Bezeichner, aber lückenhaft befüllt \\
    2 & \texttt{openfda.generic\_name:ibuprofen} & Höchste Trefferquote in der Praxis \\
    3 & \texttt{openfda.substance\_name:IBUPROFEN} & Reiner Wirkstoff, meist in Großbuchstaben gepflegt \\
    4 & \texttt{openfda.brand\_name:IBUPROFEN} & Letzte Rückfallebene \\
    \bottomrule
  \end{tabularx}
\end{table}

\paragraph{Filterung und Aggregation.} Ein zusätzlicher Existenzfilter (\texttt{\_exists\_}) auf das Feld \texttt{drug\_interactions} sorgt dafür, dass openFDA nur Fachinformationen mit Interaktionsabschnitt liefert. Ohne diesen Filter enthielten die Treffer zu Ibuprofen häufig rezeptfreie Packungsbeilagen ohne Abschnitt~7. Abgefragt werden bis zu 25 Fachinformationen, deren Absätze zusammengeführt werden, weil verschiedene Hersteller unterschiedliche Wechselwirkungen dokumentieren. Die Anfrage-URI wird mit \texttt{URLEncoder} einmalig kodiert und dem \texttt{WebClient} als fertiges \texttt{URI}-Objekt übergeben, da die Konstruktion über den \texttt{UriBuilder} zu doppelter Kodierung geführt hatte. Antworten mit Status~404 oder~400 werden als leeres Ergebnis gewertet und lösen die nächste Stufe der Kaskade aus.

\paragraph{Duplikatbereinigung.} Die Aggregation erzeugt viele nahezu identische Absätze, die sich nur in Interpunktion, Schreibweise oder einzelnen Formulierungen unterscheiden. Medcheck verwendet daher ein Fingerprint-Verfahren (Listing~\ref{lst:dedup}): Jeder Absatz wird normalisiert (Kleinschreibung, Zusammenfassen von Leerraum und Satzzeichen), die ersten 100 Zeichen dienen als Schlüssel, und pro Schlüssel bleibt der längste Absatz erhalten. Bei Ibuprofen reduzierte das Verfahren acht Absätze auf vier inhaltlich verschiedene.

\begin{lstlisting}[float=htbp, language=Java, caption={Fingerprint-basierte Duplikatbereinigung der Interaktionsabsätze.}, label={lst:dedup}]
LinkedHashMap<String, String> byFingerprint = new LinkedHashMap<>();
for (String para : paragraphs) {
    String key = fingerprint(para);        // erste 100 normalisierte Zeichen
    String existing = byFingerprint.get(key);
    if (existing == null || para.length() > existing.length()) {
        byFingerprint.put(key, para);      // längste Variante bleibt erhalten
    }
}
\end{lstlisting}

\paragraph{Caching.} Die openFDA-Abfrage ist mit bis zu vier Netzwerkaufrufen je Medikament der aufwendigste Schritt. Ihre Ergebnisse werden in der Entität \texttt{DrugInteractionText} gespeichert, mit der RxCUI (ersatzweise dem normalisierten Namen) als Schlüssel sowie kanonischem Namen, Interaktionstext, Herkunft und Abrufzeitpunkt. Da aggregierte Texte mehr als 30\,000 Zeichen umfassen können, ist das Textfeld als Large Object (\texttt{@Lob}) deklariert. Eine frühere Version speicherte Ergebnisse paarweise. Der Cache pro Medikament benötigt bei \(n\) Medikamenten höchstens \(n\) statt \(n(n-1)/2\) Einträge, und jeder Eintrag ist in beliebigen Kombinationen wiederverwendbar. Die paarweise Bewertung erfolgt ausschließlich im Sprachmodell.

\section{Anbieterunabhängige LLM-Schicht}

Die Anforderung NF2 ist mit dem Strategie-Entwurfsmuster umgesetzt. Das Interface \texttt{LlmProvider} (Listing~\ref{lst:provider}) definiert drei Operationen, und der \texttt{InteractionService} kennt ausschließlich dieses Interface. Die Klassen \texttt{AnthropicService}, \texttt{GeminiService} und \texttt{OllamaService} implementieren es. Über die Annotation \texttt{@ConditionalOnProperty(name = "llm.provider", ...)} wird beim Start genau eine Implementierung registriert; ein Anbieterwechsel ist damit eine reine Konfigurationsänderung. Tabelle~\ref{tab:anbieter} zeigt, wie stark sich die Anbieter im Transportformat unterscheiden. Alle drei verwenden dennoch identische Prompts und dieselbe Auswertungslogik.

\begin{lstlisting}[float=htbp, language=Java, caption={Das Interface \texttt{LlmProvider}.}, label={lst:provider}]
public interface LlmProvider {
    boolean isConfigured();
    LlmVerdict explainInteraction(List<String> drugNames,
                                  List<String> drugTexts,
                                  SupportedLanguage language);
    String summarizeSideEffects(String drugName, String labelText,
                                SupportedLanguage language);
}
\end{lstlisting}

\begin{table}[htbp]
  \centering
  \small
  \caption{Vergleich der implementierten LLM-Anbieter.}
  \label{tab:anbieter}
  \begin{tabularx}{\textwidth}{@{}l L L L@{}}
    \toprule
    Merkmal & Anthropic Claude & Google Gemini & Ollama \\
    \midrule
    Endpunkt          & \texttt{/v1/messages} & \texttt{/v1beta/models/}\allowbreak\texttt{\{m\}:generateContent} & \texttt{/api/chat} \\
    Authentifizierung & Header \texttt{x-api-key} & Query-Parameter \texttt{key} & keine \\
    System-Prompt     & Feld \texttt{system} & Objekt \texttt{systemInstruction} & Nachricht mit Rolle \texttt{system} \\
    Ausführung        & Cloud & Cloud & lokal \\
    Timeout           & 15\,s & 15\,s & 180\,s \\
    \bottomrule
  \end{tabularx}
\end{table}

Die lokale Variante über Ollama ist im medizinischen Kontext besonders interessant, weil keine Daten den Rechner verlassen. Erkauft wird dies mit deutlich längeren Antwortzeiten und einer geringeren Formattreue kleiner Modelle (Abschnitt~\ref{sec:herausforderungen}).

\section{Prompt Engineering}

Alle Prompts sind in der Klasse \texttt{AnthropicPrompts} gebündelt, damit sie als eigenständiges Artefakt versioniert und dokumentiert werden können. Der Entwurf folgt fünf Prinzipien. (1)~\emph{Grounding}: Das Modell darf ausschließlich den übergebenen openFDA-Text auswerten. (2)~\emph{Konservative Einstufung}: Bei mehrdeutiger Evidenz ist die höhere Stufe zu wählen, weil eine falsch-negative Einstufung in diesem Anwendungsfeld schwerer wiegt als eine falsch-positive. (3)~\emph{Strukturierte Ausgabe}: Die Bewertung wird als JSON-Objekt mit den Schlüsseln \texttt{severity} und \texttt{summary} angefordert. (4)~\emph{Englische Instruktion bei wählbarer Ausgabesprache}: Die Anweisungen sind englisch formuliert, weil Sprachmodelle englischen Instruktionen am zuverlässigsten folgen; die Zielsprache wird über den Platzhalter \texttt{\{LANGUAGE\}} eingesetzt und zweimal genannt. (5)~\emph{Verweis auf Fachpersonal}: Jede Antwort fordert zur Rücksprache mit Ärztin, Arzt oder Apotheke auf. Listing~\ref{lst:prompt} zeigt einen Auszug.

\begin{lstlisting}[float=htbp, caption={Auszug aus dem System-Prompt für die Bewertung einer Kombination.}, label={lst:prompt}]
Analyse ONLY the text provided; do not invent facts, do not consult
external knowledge, and do not give personal medical advice.
[...]
  2. Write a short explanation (2-4 sentences) for a layperson.
Rules:
  - Answer STRICTLY in {LANGUAGE}, in plain, understandable language.
  - Be CONSERVATIVE: if the evidence is ambiguous, err on the
    higher severity. Never downgrade a clear warning.
  - Always tell the reader to consult a doctor or pharmacist for
    medical decisions.
Output format - return ONLY one JSON object with exactly these two
keys, no markdown fences, no comments, no extra prose:
  { "severity": "LOW" | "MEDIUM" | "HIGH",
    "summary": "..." }
The value of "summary" must be written in {LANGUAGE}.
\end{lstlisting}

Ein zweiter Prompt erzeugt je Medikament eine Zusammenfassung der wichtigsten Warnhinweise in zwei bis vier Sätzen. Vor dem Aufruf wird der openFDA-Text auf 3\,000 Zeichen je Medikament gekürzt. Das senkt Kosten und Antwortzeit; das Frontend erhält weiterhin den vollständigen Text.

\section{Auswertung der Modellantwort und kontrollierte Degradation}

Auch bei eindeutiger Anweisung halten sich Sprachmodelle nicht immer an das geforderte Format. Typisch sind einleitende Sätze vor dem JSON-Objekt oder, bei kleinen lokalen Modellen, frei erfundene JSON-Strukturen. Die Hilfsklasse \texttt{LlmParsing} ist daher defensiv aufgebaut: Ein regulärer Ausdruck extrahiert das erste Objekt zwischen geschweiften Klammern, Jackson liest daraus die beiden Felder, und \texttt{Severity.fromString} bildet gängige Synonyme ab (etwa \texttt{MODERATE} auf \texttt{MEDIUM}). Jeder Fehler führt zum Ergebnis \texttt{UNKNOWN}. Dieser Wert ist kein Risikourteil, sondern kennzeichnet das Fehlen einer verwertbaren Einstufung und wird im Frontend bewusst nicht als geringes Risiko dargestellt.

Für \(n\) Medikamente sind \(n+1\) Modellaufrufe nötig. Da sie unabhängig voneinander sind, startet der \texttt{InteractionService} sie mit \texttt{CompletableFuture} parallel und wartet auf alle Ergebnisse. Schlägt ein Aufruf fehl oder ist kein Anbieter konfiguriert, liefert der Service eine Antwort mit \texttt{severity = UNKNOWN}, ohne Zusammenfassungen und mit \texttt{llmAvailable = false}, aber mit allen bereits beschafften openFDA-Texten. Der Endpunkt antwortet in diesem Fall weiterhin mit Status~200.

\section{Mehrsprachigkeit, Validierung und Fehlerbehandlung}

Die Ausgabesprachen sind im Enum \texttt{SupportedLanguage} modelliert. Jeder Eintrag enthält den ISO-639-1-Code, die Eigenbezeichnung für die Oberfläche und den englischen Namen für den Prompt, beispielsweise \texttt{TURKISH("tr", "Türkçe", "Turkish")}. Hinterlegt sind Deutsch, Englisch, Spanisch, Französisch, Italienisch, Portugiesisch, Türkisch, Niederländisch, Polnisch, Russisch und Arabisch. Gegenüber einer Konfigurationsdatei bietet das Enum Typsicherheit; da der englische Name unmittelbar in den Prompt eingeht, ist er Teil der fachlichen Logik.

Eingaben prüft die Klasse \texttt{DrugValidator}, bevor sie die Geschäftslogik erreichen: Namen müssen 2 bis 100 Zeichen lang sein und dürfen nur Zeichen einer Positivliste enthalten, die neben Buchstaben und Ziffern auch die in RxNorm-Bezeichnungen üblichen Klammern, Schrägstriche und Prozentzeichen umfasst. Eine Liste muss 2 bis 20 Einträge enthalten. Ein zentraler \texttt{GlobalExceptionHandler} übersetzt Ausnahmen in einheitliche Fehlerantworten: Validierungsfehler in Status~400 mit Detailliste, Fehler von RxNorm oder openFDA in Status~502 und unerwartete Fehler in Status~500 ohne interne Details.

\section{Frontend}

Das Frontend ist als Angular-Anwendung mit Standalone Components aufgebaut und gliedert sich in \texttt{core} (HTTP-Services), \texttt{features} (Such- und Ergebnisansicht) und \texttt{shared} (TypeScript-Interfaces, die die Backend-DTOs abbilden). Auf eine Komponentenbibliothek wurde nach einer ersten Version verzichtet; die Bedienelemente sind mit nativem HTML und eigenem CSS umgesetzt, was die Bundle-Größe reduziert und volle Kontrolle über Darstellung und Verhalten gibt.

Die Suchansicht beginnt mit zwei Eingabefeldern, weitere lassen sich hinzufügen. Jedes Feld besitzt eine eigene Autovervollständigung als reaktiven RxJS-Datenstrom: Eingaben werden verzögert, erst ab drei Zeichen an das Backend gesendet, und \texttt{switchMap} verwirft veraltete Anfragen. Die Vorschlagsliste ist vollständig per Tastatur bedienbar. Darunter wählt eine Reihe von Schaltflächen die Antwortsprache, wobei jede Sprache in ihrer Eigenbezeichnung erscheint.

Die Ergebnisansicht zeigt das Wichtigste ohne weitere Interaktion: eine Zusammenfassungskarte mit Ampel, textueller Risikobezeichnung und der Erklärung des Modells (Tabelle~\ref{tab:ampel}). Die Einstufung wird nie allein über Farbe vermittelt, sondern zusätzlich über die Position der Leuchte und eine Textbezeichnung. Darunter erscheint für jedes Medikament ein aufklappbares Element, in dem ein Umschalter zwischen \enquote{AI Kurzfassung} und \enquote{FDA-Original} wechselt. Im Degradationsfall zeigt die Karte den Hinweis \enquote{KI derzeit nicht verfügbar}, während die Originaltexte zugänglich bleiben.

\begin{table}[htbp]
  \centering
  \small
  \caption{Abbildung der Risikostufen auf die Ampeldarstellung.}
  \label{tab:ampel}
  \begin{tabular}{@{}llll@{}}
    \toprule
    Stufe & Ampel & Bezeichnung & Kartenhintergrund \\
    \midrule
    \texttt{HIGH}    & \riskdot{riskhigh}\ oben, rot     & Hohes Risiko     & hellrot \\
    \texttt{MEDIUM}  & \riskdot{riskmid}\ Mitte, gelb-orange& Mittleres Risiko & hellorange \\
    \texttt{LOW}     & \riskdot{risklow}\ unten, grün    & Niedriges Risiko & hellgrün \\
    \texttt{UNKNOWN} & \riskoff\ alle Leuchten aus       & Risiko unbekannt & neutral grau \\
    \bottomrule
  \end{tabular}
\end{table}

\section{Qualitätssicherung}

Die Anbindung externer Dienste ist der fehleranfälligste Teil des Systems und wurde am gründlichsten getestet. Alle Service-Tests verwenden den OkHttp \texttt{MockWebServer}, einen lokal gestarteten HTTP-Server mit vorgegebenen Antworten. So lassen sich Statuscodes, Verzögerungen und fehlerhafte Modellantworten reproduzierbar prüfen, ohne Netzwerkzugriff oder API-Schlüssel. Insgesamt umfasst die Testsuite 31 Fälle: vier für RxNorm, sieben für openFDA, acht für Anthropic sowie je sechs für Gemini und Ollama. Besonderes Gewicht liegt auf den Negativfällen: Für jeden Anbieter wird geprüft, dass ein ungültiges Ausgabeformat zur Einstufung \texttt{UNKNOWN} führt und dass ein Timeout als anbieterspezifische Ausnahme erscheint, die der \texttt{InteractionService} abfangen kann.

\section{Herausforderungen bei der Umsetzung}
\label{sec:herausforderungen}

Ein erheblicher Teil des Aufwands entfiel auf das Verhalten der externen Dienste im Realbetrieb. Tabelle~\ref{tab:probleme} fasst die wichtigsten Probleme zusammen. Mehrere davon traten erst mit echten Daten auf und wären mit reinen Testattrappen nicht entdeckt worden.

\begin{table}[htbp]
  \centering
  \small
  \caption{Aufgetretene Probleme, Ursachen und Lösungen.}
  \label{tab:probleme}
  \begin{tabularx}{\textwidth}{@{}L L L@{}}
    \toprule
    Problem & Ursache & Lösung \\
    \midrule
    openFDA ohne Treffer trotz vorhandener Daten & Doppelte URL-Kodierung & Explizite Kodierung, Übergabe als \texttt{URI} \\
    Status 400 statt 404 & Eigenheit von openFDA & Beide Status als leeres Ergebnis werten \\
    Suche über RxCUI erfolglos & Feld lückenhaft befüllt & Vierstufige Abfragekaskade \\
    Treffer ohne Interaktionsabschnitt & Rezeptfreie Packungsbeilagen & Filter \texttt{\_exists\_} \\
    Viele nahezu identische Absätze & Aggregation über 25 Labels & Fingerprint-Deduplizierung \\
    Abbruch beim Lesen großer Antworten & Pufferlimit des \texttt{WebClient} & Puffer auf 16\,MB erhöht \\
    Cache-Einträge nicht gespeichert & Spalte mit 8\,000 Zeichen zu kurz & Textfeld als \texttt{@Lob} \\
    Cloud-Modell lehnt Anfrage ab (404, 429) & Modell abgekündigt, kein Kontingent & Modell konfigurierbar, Anbieter austauschbar \\
    Lokales Modell überschreitet Timeout & CPU-Inferenz mit langem Kontext & Eingabekürzung, längerer Timeout, kleineres Modell \\
    Kleines Modell erfindet eigenes JSON-Schema & Geringe Formattreue & Defensive Auswertung, Rückfall auf \texttt{UNKNOWN} \\
    \bottomrule
  \end{tabularx}
\end{table}

Daraus lassen sich zwei allgemeine Erkenntnisse ableiten. Erstens ist die Dokumentation öffentlicher Schnittstellen keine vollständige Spezifikation; die Protokollierung der tatsächlich gesendeten Anfragen war der schnellste Weg zur Fehlerursache. Zweitens ist die Ausgabe eines Sprachmodells als nicht vertrauenswürdige Eingabe zu behandeln: Jede Annahme über ihr Format muss abgesichert sein, und ein Formatfehler darf nie die gesamte Anfrage scheitern lassen.

\chapter{Evaluation}
\label{ch:evaluation}

Dieses Kapitel beantwortet FF2 anhand der in Abschnitt~\ref{sec:befragungsdesign} beschriebenen Befragung. Ergänzend wird der Prototyp gegen die Anforderungen aus Kapitel~\ref{ch:methodik} bewertet.

\section{Stichprobe}

An der Befragung nahmen elf Personen teil, die im Folgenden in der Reihenfolge der Befragung als K1 bis K11 bezeichnet werden. Acht Personen gehören der Altersgruppe 18 bis 29 Jahre an, zwei der Gruppe 45 bis 59 Jahre und eine der Gruppe 16 bis 18 Jahre. Nur zwei Personen nehmen regelmäßig Medikamente ein, sechs haben bereits online nach Medikamenteninformationen gesucht. Ihre Sicherheit im Umgang mit medizinischen Informationen schätzten die Teilnehmenden im Mittel mit 3{,}55 ein (Median~3). Variante~A wurde acht Personen zuerst gezeigt, Variante~B drei Personen.

\section{Bewertung der Aussagen}

Tabelle~\ref{tab:itemergebnisse} zeigt die Bewertungen der acht Aussagen. Unter Berücksichtigung der negativ formulierten Aussage V3 schneidet Variante~B bei sieben von acht Aussagen besser ab. Am deutlichsten sind die Unterschiede bei der Verständlichkeit: B wurde als leichter verständlich (V1) und begrifflich klarer (V2) bewertet und musste seltener mehrfach gelesen werden (V3). Ebenfalls deutlich ist der Unterschied bei N2: Nach dem Lesen von B wussten die Teilnehmenden besser, was als Nächstes zu tun ist. Bei den Vertrauensaussagen T1 und T2 liegen beide Varianten nahe beieinander. Die einzige Aussage, bei der A besser abschneidet, ist N1: Variante~A half etwas mehr, das Risiko einzuschätzen.

\begin{table}[htbp]
  \centering
  \small
  \caption{Mittelwert (M), Standardabweichung (SD) und Median (Md) je Aussage auf der Skala von 1 bis 5. Bei V3 ist ein niedriger Wert positiv.}
  \label{tab:itemergebnisse}
  \begin{tabular}{@{}llcccc@{}}
    \toprule
    & & \multicolumn{2}{c}{Variante A} & \multicolumn{2}{c}{Variante B} \\
    \cmidrule(lr){3-4}\cmidrule(l){5-6}
    Kürzel & Aussage (Kurzform) & M (SD) & Md & M (SD) & Md \\
    \midrule
    V1 & leicht zu verstehen          & 3{,}73 (1{,}19) & 4 & 4{,}45 (0{,}93) & 5 \\
    V2 & Begriffe klar                & 3{,}73 (1{,}01) & 4 & 4{,}36 (0{,}92) & 5 \\
    V3 & mehrfach lesen (negativ)     & 2{,}55 (1{,}21) & 2 & 1{,}91 (1{,}22) & 1 \\
    T1 & vertraue der Richtigkeit     & 3{,}91 (1{,}14) & 4 & 4{,}09 (1{,}04) & 4 \\
    T2 & seriös und glaubwürdig       & 4{,}09 (1{,}04) & 4 & 4{,}27 (1{,}01) & 5 \\
    N1 & Risiko einschätzen           & 4{,}18 (1{,}25) & 5 & 4{,}09 (1{,}04) & 4 \\
    N2 & weiß, was zu tun ist         & 3{,}82 (1{,}25) & 4 & 4{,}36 (1{,}03) & 5 \\
    N3 & alle Informationen enthalten & 3{,}73 (1{,}35) & 4 & 4{,}00 (1{,}26) & 4 \\
    \bottomrule
  \end{tabular}
\end{table}

Abbildung~\ref{fig:dimensionen} fasst die Aussagen zu den drei Dimensionen und einem Gesamtwert zusammen. Variante~B liegt in allen Dimensionen vorn, am deutlichsten bei der Verständlichkeit (4{,}30 gegenüber 3{,}64). Tabelle~\ref{tab:wilcoxon} zeigt die Ergebnisse des Wilcoxon-Vorzeichen-Rang-Tests. Keiner der Unterschiede ist auf dem Niveau \(\alpha = 0{,}05\) signifikant. Die Effektstärke liegt für Verständlichkeit (\(r = 0{,}47\)) und Gesamtwert (\(r = 0{,}40\)) im mittleren Bereich, für Vertrauen und Nützlichkeit im kleinen Bereich. Bei elf Personen lässt sich ein Effekt dieser Größe statistisch nicht absichern. Die Ergebnisse sind daher als Hinweise zu verstehen, nicht als Nachweis.

\begin{figure}[htbp]
  \centering
  \begin{tikzpicture}
    \begin{axis}[
        ybar, bar width=17pt,
        width=0.86\textwidth, height=4.8cm,
        ymin=1, ymax=5.3, ytick={1,2,3,4,5},
        ylabel={Mittelwert (1--5)}, ylabel style={font=\small},
        xtick={1,2,3,4}, xmin=0.5, xmax=4.5,
        xticklabels={Verständlichkeit, Vertrauen, Nützlichkeit, Gesamt},
        tick label style={font=\small}, xtick style={draw=none},
        ymajorgrids, grid style={black!12}, axis lines*=left,
        yticklabel={\pgfmathprintnumber[use comma,fixed,precision=0]{\tick}},
        nodes near coords={\pgfmathprintnumber[use comma,fixed,fixed zerofill,precision=2]{\pgfplotspointmeta}},
        every node near coord/.append style={font=\scriptsize},
        legend style={at={(0.5,1.03)}, anchor=south, legend columns=2, draw=none, font=\small,
                      /tikz/every even column/.append style={column sep=1.2em}},
        area legend
      ]
      \addplot[fill=varA, draw=none] coordinates {(1,3.64) (2,4.00) (3,3.91) (4,3.83)};
      \addplot[fill=varB, draw=none] coordinates {(1,4.30) (2,4.18) (3,4.15) (4,4.22)};
      \legend{Variante A (Drugs.com), Variante B (Medcheck)}
    \end{axis}
  \end{tikzpicture}
  \caption{Mittelwerte der Dimensionen und des Gesamtwerts. V3 ist umkodiert.}
  \label{fig:dimensionen}
\end{figure}

\begin{table}[htbp]
  \centering
  \small
  \caption{Wilcoxon-Vorzeichen-Rang-Test (exakt, zweiseitig). \(n\): Personen mit von null verschiedener Differenz; \(W^{+}\), \(W^{-}\): Rangsummen der Differenzen zugunsten von B bzw.\ A; \(r\): Rang-biseriale Korrelation.}
  \label{tab:wilcoxon}
  \begin{tabular}{@{}lcccccc@{}}
    \toprule
    Dimension & A: M (SD) & B: M (SD) & \(n\) & \(W^{+}\) / \(W^{-}\) & \(p\) & \(r\) \\
    \midrule
    Verständlichkeit & 3{,}64 (0{,}94) & 4{,}30 (0{,}85) & 10 & 40{,}5 / 14{,}5 & 0{,}203 & 0{,}47 \\
    Vertrauen        & 4{,}00 (1{,}00) & 4{,}18 (1{,}01) & 7  & 18{,}0 / 10{,}0 & 0{,}531 & 0{,}29 \\
    Nützlichkeit     & 3{,}91 (1{,}16) & 4{,}15 (1{,}07) & 11 & 40{,}5 / 25{,}5 & 0{,}526 & 0{,}23 \\
    Gesamt           & 3{,}83 (0{,}86) & 4{,}22 (0{,}86) & 10 & 38{,}5 / 16{,}5 & 0{,}287 & 0{,}40 \\
    \bottomrule
  \end{tabular}
\end{table}

Die Mittelwerte verdecken eine große Streuung zwischen den Personen (Abbildung~\ref{fig:differenzen}). Sieben Personen bewerteten B im Gesamtwert besser, drei schlechter und eine gleich. Die Differenzen reichen von rund \(-2{,}1\) Punkten (K3) bis \(+1{,}9\) Punkten (K5).

\begin{figure}[htbp]
  \centering
  \begin{tikzpicture}
    \begin{axis}[
        ybar, bar width=13pt, bar shift=0pt,
        width=0.86\textwidth, height=4.5cm,
        ymin=-3.0, ymax=2.6, ytick={-2,-1,0,1,2},
        ylabel={Differenz B $-$ A}, ylabel style={font=\small},
        xtick={1,...,11}, xmin=0.4, xmax=11.6,
        xticklabels={K1,K2,K3,K4,K5,K6,K7,K8,K9,K10,K11},
        tick label style={font=\small}, xtick style={draw=none},
        ymajorgrids, grid style={black!12}, axis lines*=left,
        extra y ticks={0}, extra y tick labels={},
        extra y tick style={grid style={black!60, line width=0.5pt}},
        yticklabel={\pgfmathprintnumber[use comma,fixed,precision=0]{\tick}},
        nodes near coords={\pgfmathprintnumber[use comma,fixed,fixed zerofill,precision=1]{\pgfplotspointmeta}},
        every node near coord/.append style={font=\scriptsize},
        legend style={at={(0.5,1.03)}, anchor=south, legend columns=2, draw=none, font=\small,
                      /tikz/every even column/.append style={column sep=1.2em}},
        area legend
      ]
      \addplot[fill=orderA, draw=none] coordinates
        {(1,1.5) (3,-2.125) (5,1.875) (7,1.625) (8,0.125) (9,0.75) (10,0.375) (11,1.125)};
      \addplot[fill=orderB, draw=none] coordinates {(2,-0.875) (4,0) (6,-0.125)};
      \legend{Variante A zuerst gelesen, Variante B zuerst gelesen}
    \end{axis}
  \end{tikzpicture}
  \caption{Differenz der Gesamtwerte (B $-$ A) je Person. Positive Werte bedeuten eine bessere Bewertung von Variante~B.}
  \label{fig:differenzen}
\end{figure}

\section{Direkter Vergleich und Reihenfolgeeffekt}

Im direkten Vergleich fanden acht von elf Personen Variante~B verständlicher (Tabelle~\ref{tab:vergleich}). Beim Vertrauen kehrt sich das Bild knapp um: Sechs Personen würden eher Variante~A vertrauen. Acht Personen entschieden sich in allen vier Fragen für dieselbe Variante, fünf für B und drei für A. Die übrigen drei (K7, K8, K10) fanden B verständlicher, würden aber eher A vertrauen. Verständlichkeit und Vertrauen fallen also bei einem Teil der Personen auseinander. Keine der Präferenzen weicht signifikant von einer Gleichverteilung ab.

\begin{table}[htbp]
  \centering
  \small
  \caption{Präferenzen im direkten Vergleich mit zweiseitigem Binomialtest gegen \(p_0 = 0{,}5\).}
  \label{tab:vergleich}
  \begin{tabular}{@{}lccc@{}}
    \toprule
    Frage & A & B & \(p\) \\
    \midrule
    Welche Erklärung fanden Sie verständlicher?                  & 3 & 8 & 0{,}227 \\
    Welcher Erklärung würden Sie eher vertrauen?                 & 6 & 5 & 1{,}000 \\
    Welche half Ihnen mehr, eine Entscheidung zu treffen?        & 4 & 7 & 0{,}549 \\
    Welche würden Sie einem Freund/Familienmitglied empfehlen?   & 4 & 7 & 0{,}549 \\
    \bottomrule
  \end{tabular}
\end{table}

Da die Reihenfolge mit 8:3 nicht ausbalanciert war, ist ein möglicher Reihenfolgeeffekt gesondert zu betrachten. Acht von elf Personen empfahlen die Variante, die sie als zweite gelesen hatten (\(p = 0{,}227\)). In der Gruppe, die A zuerst las, bewerteten sieben von acht Personen B im Gesamtwert besser, im Mittel um 0{,}66 Punkte. Sechs von ihnen empfahlen B. In der Gruppe, die B zuerst las, bewertete keine Person B besser, im Mittel lag B um 0{,}33 Punkte darunter, und zwei der drei Personen empfahlen A. Mit drei Personen ist diese Gruppe zu klein für belastbare Aussagen. Das Muster ist jedoch mit einer inhaltlichen Erklärung vereinbar: Wer A zuerst gelesen hat, kennt die Details bereits und erlebt B als klare Zusammenfassung. Wer mit B beginnt, vermisst die Details und findet sie anschließend in A.

\section{Verständnistest und offene Antworten}

Die offenen Fragen des Verständnistests zeigen, welche Botschaft bei den Lesenden ankam. Offensichtliche Tippfehler in den folgenden Zitaten wurden korrigiert. Als nächsten Schritt nannten acht Personen die Rücksprache mit Ärztin, Arzt oder Apotheke (K1, K2, K4, K5, K6, K7, K10, K11). K1 und K4 entnahmen Variante~B zudem ausdrücklich, die Medikamente nicht zu kombinieren. Variante~A hinterließ dagegen bei einigen Personen Unsicherheit. K1 antwortete \enquote{Ich weiß es nicht wirklich, zum Arzt?}, und K9 empfand die Formulierung \enquote{may increase} als unklar. K4 entnahm A die Möglichkeit, nach ärztlicher Rücksprache die Dosierung anzupassen, und schätzte die Kombination danach als \enquote{nicht so schlimm} ein, nach B dagegen als \enquote{wahrscheinlich schlimm}.

Bei der Einschätzung der Gefahr ergibt sich ein gemischtes Bild. K1, K4, K5 und K9 leiteten aus Variante~B eine klarere oder höhere Gefahr ab. K9 schrieb, die deutliche Sprache halte von einer Kombination ab (\enquote{strong language pushes me away from it}). Umgekehrt begründeten K2, K3 und K6 ihre Einschätzung mit Details, die nur Variante~A enthält, nämlich den beschriebenen Symptomen und dem Hinweis auf den Magen-Darm-Trakt. K2 sagte zu Variante~B: \enquote{Da würde ich mir jetzt keine Gedanken machen}, weil der Text nicht vermittle, dass die Kombination richtig gefährlich sei.

Auffällig ist die Begründung von K7, die Kombination sei \enquote{sehr gefährlich, da beide Wirkstoffe zur Blutgerinnung führen können}. Tatsächlich beeinträchtigen beide Wirkstoffe die Blutstillung, Warfarin über die Gerinnungsfaktoren und Ibuprofen über die Funktion der Blutplättchen. Variante~B formuliert lediglich, dass beide die Blutgerinnung \enquote{beeinflussen}, und lässt die Wirkrichtung damit offen. Die Antworten von K8 zur ersten Frage sind nicht auswertbar, weil die Frage als Bitte um Verbesserungsvorschläge verstanden wurde. Bei K11 sind die Bezeichnungen A und B in Teil~4 und~5 offensichtlich vertauscht, denn die zitierten Formulierungen stammen jeweils aus der anderen Variante. Diese Antworten wurden mit getauschten Bezeichnungen ausgewertet.

Die Antworten auf die offenen Fragen aus Teil~5 und die Begründungen im direkten Vergleich lassen sich zu vier Themen zusammenfassen (Tabelle~\ref{tab:themen}). Die Themen sind nicht widersprüchlich, sondern beschreiben einen Zielkonflikt: Dieselbe Kürze, die B verständlich macht, führt dazu, dass Details fehlen. Mehrere Personen nannten beide Seiten. So beschrieb K9 Variante~B als \enquote{short and very clear}, vermisste darin aber die bei A genannten Symptome.

\begin{table}[htbp]
  \centering
  \small
  \caption{Themen der offenen Antworten.}
  \label{tab:themen}
  \begin{tabularx}{\textwidth}{@{}>{\raggedright\arraybackslash}p{4.1cm} >{\raggedright\arraybackslash}p{2.6cm} L@{}}
    \toprule
    Thema & Personen & Beispiel \\
    \midrule
    B ist kurz, klar und ohne Fachbegriffe & K1, K4, K5, K7, K9, K10, K11 & \enquote{Es war kurz und knapp und die Infos wurden übermittelt.} (K5) \\
    B fehlen Details und konkrete Warnzeichen & K2, K3, K6, K8, K9, K11 & \enquote{Es war sehr komprimiert, aber sehr oberflächlich erklärt.} (K3) \\
    A ist zu lang oder zu fachsprachlich & K1, K4, K5, K8, K9, K10 & \enquote{Mit jedem zusätzlichen Satz hatte ich mehr Fragen als eine Erklärung.} (K5) \\
    Details schaffen Vertrauen und Risikobewusstsein & K2, K3, K6, K8, K11 & \enquote{Variant A is more complicated to understand but it makes the product more trustworthy.} (K8) \\
    \bottomrule
  \end{tabularx}
\end{table}

Als mögliche Nutzungssituationen nannten die Teilnehmenden das Kombinieren von Medikamenten zu Hause (K1, K2), die Einnahme von Schmerzmitteln zusammen mit anderen Medikamenten (K7), die Vorbereitung eines Apothekenbesuchs (K8), Nahrungsergänzungsmittel (K9), die Einnahme von mehr als zwei Medikamenten (K10) und die Prüfung eines ärztlichen Rezepts (K11). K9 würde das Werkzeug auch den eigenen Eltern anstelle des Beipackzettels empfehlen. K3 machte die Nutzung davon abhängig, ob das Werkzeug vertrauenswürdig ist. K4 wurde zusätzlich die Oberfläche des Prototyps gezeigt. Zur Ampel äußerte K4: \enquote{Das macht es auf jeden Fall einfacher. Vor allem für Menschen, die sich nicht so gut auskennen.}

\section{Diskussion}

\paragraph{Verständlichkeit.} Die Ergebnisse deuten darauf hin, dass Medcheck das Ziel einer laienverständlichen Erklärung erreicht. Der Vorteil von Variante~B zeigt sich in allen drei Verständlichkeitsaussagen, im direkten Vergleich und in den offenen Antworten. Das stützt die Einschätzung von Mohammed et al., dass LLMs als ergänzende Erklärwerkzeuge dienen können~\cite{MOHAMMED2026100733}, ebenso den Ansatz von Almeida et al., Arzneimittelinformationen mit generativer KI für Patientinnen und Patienten aufzubereiten~\cite{Almeida2025}. Statistisch abgesichert ist dieser Vorteil bei der vorliegenden Stichprobe allerdings nicht.

\paragraph{Vertrauen.} Beim Vertrauen erreicht die kurze Erklärung nicht mehr als der etablierte Text. Mehrere Teilnehmende verbanden Vertrauen mit Ausführlichkeit, am deutlichsten K8. Zu berücksichtigen ist, dass in der Befragung nur der Text gezeigt wurde. Dass Variante~B auf der amtlichen Fachinformation beruht und diese in der Anwendung mit einem Klick einsehbar ist, war für die Teilnehmenden nicht erkennbar.

\paragraph{Fehlende Warnzeichen als Folge des Groundings.} Der häufigste Kritikpunkt, die fehlenden Warnzeichen, ist eine direkte Folge des Entwurfs. Medcheck wertet nur den Interaktionsabschnitt der Fachinformation aus, und das Modell ist angewiesen, nichts hinzuzufügen. Warnzeichen einer Blutung stehen in Fachinformationen jedoch in der Regel nicht im Interaktionsabschnitt, sondern bei den Warnhinweisen und in den Hinweisen zur Patientenberatung. Variante~A ist dagegen ein für Patientinnen und Patienten verfasster Text, der diese Warnzeichen ausdrücklich nennt. Das Grounding verhindert also erfundene Inhalte, begrenzt die Antwort aber auf das, was der ausgewertete Abschnitt enthält. Die Lösung liegt daher nicht in einem freieren Prompt, sondern in einer breiteren Datengrundlage.

\paragraph{Deutlichkeit der Warnung.} Die deutliche Formulierung von Variante~B wurde von mehreren Personen als klare Warnung verstanden (K1, K4, K5, K9). Die vorsichtigere Formulierung von Variante~A (\enquote{may increase}) hinterließ zusammen mit den genannten Handlungsoptionen bei einigen Personen den Eindruck eines beherrschbaren Risikos. Für eine Kombination mit bekanntermaßen erhöhtem Blutungsrisiko ist die Wirkung von B erwünscht. Die Literatur zur Alert Fatigue zeigt jedoch, dass zu viele oder zu drastische Warnungen die Akzeptanz untergraben~\cite{VillaZapata}, und frei antwortende Modelle neigen zur Überschätzung von Wechselwirkungen~\cite{MOHAMMED2026100733}. Ob die konservative Einstufungsregel des Prompts bei harmlosen Kombinationen zu unnötigen Warnungen führt, wurde in dieser Arbeit nicht untersucht.

\paragraph{Geschichtete Darstellung.} Insgesamt haben die beiden Varianten komplementäre Stärken. B ist verständlicher und gibt eine klare Handlungsrichtung vor. A vermittelt durch Details und Warnzeichen Vertrauen und erleichtert die Einschätzung des Risikos. Der beobachtete Reihenfolgeeffekt weist in dieselbe Richtung, denn B wurde vor allem dann besser bewertet, wenn die Details aus A bereits bekannt waren. Dies spricht für eine geschichtete Darstellung: Zuerst erscheint die kurze Zusammenfassung, Warnzeichen und Originaltext sind unmittelbar darunter verfügbar. Die Oberfläche von Medcheck ist bereits so aufgebaut, bewertet wurde in der Befragung jedoch nur der Text der Zusammenfassung.

\section{Technische Bewertung}

Tabelle~\ref{tab:erfuellung} bewertet den Prototyp gegen die Anforderungen aus Tabelle~\ref{tab:anforderungen}. Alle funktionalen Anforderungen sind umgesetzt. Bei den nichtfunktionalen Anforderungen bestehen zwei Einschränkungen: Die kontrollierte Degradation ist auf Ebene der einzelnen Anbieter automatisiert getestet, im \texttt{InteractionService} jedoch nicht. Das Grounding ist im Prompt vorgegeben, seine Einhaltung wurde aber nicht systematisch geprüft.

\begin{table}[htbp]
  \centering
  \small
  \caption{Erfüllung der Anforderungen.}
  \label{tab:erfuellung}
  \begin{tabularx}{\textwidth}{@{}l l L@{}}
    \toprule
    ID & Status & Nachweis \\
    \midrule
    F1  & erfüllt    & RxNorm-Anbindung mit Autovervollständigung, vier Tests \\
    F2  & erfüllt    & Dynamische Eingabefelder, \texttt{DrugValidator} mit 2 bis 20 Einträgen \\
    F3  & erfüllt    & Klassifikation \texttt{LOW}/\texttt{MEDIUM}/\texttt{HIGH}, Ampel, \texttt{UNKNOWN} bei Formatfehlern \\
    F4  & erfüllt    & Elf Ausgabesprachen über \texttt{SupportedLanguage} \\
    F5  & erfüllt    & Umschaltbare Ansicht \enquote{FDA-Original} \\
    NF1 & weitgehend & Timeouts und Formatfehler je Anbieter getestet; Degradation im \texttt{InteractionService} nicht automatisiert getestet \\
    NF2 & erfüllt    & Drei Anbieter, Wechsel über \texttt{llm.provider} ohne Codeänderung \\
    NF3 & teilweise  & Grounding im Prompt vorgegeben; Einhaltung nicht systematisch geprüft \\
    NF4 & erfüllt    & Cache pro Medikament, Kürzung der Modelleingabe, lokales Modell möglich \\
    \bottomrule
  \end{tabularx}
\end{table}

Systematische Messungen zu Antwortzeiten und Formattreue der Anbieter wurden nicht durchgeführt. Qualitativ zeigte sich während der Entwicklung ein deutlicher Unterschied zwischen Cloud- und lokalem Betrieb. Probleme bei der Cloud-Anbindung waren vor allem organisatorischer Natur, etwa abgekündigte Modellnamen und fehlendes Kontingent. Beim lokalen Modell betrafen sie dagegen Antwortzeit und Formattreue (Tabelle~\ref{tab:probleme}). Die lokale Variante ist aus Sicht des Datenschutzes attraktiv, setzt für einen praxistauglichen Einsatz aber geeignete Hardware voraus.

\section{Limitationen}

Die Ergebnisse der Befragung unterliegen mehreren Einschränkungen:

\begin{itemize}
  \item \textbf{Stichprobe.} Elf Personen aus dem persönlichen Umfeld, überwiegend 18 bis 29 Jahre alt und ohne regelmäßige Medikation, bilden die eigentliche Risikogruppe der Polypharmazie nicht ab. Die geringe Fallzahl erlaubt nur den Nachweis großer Effekte.
  \item \textbf{Untersuchungsgegenstand.} Bewertet wurden ein einziges Arzneimittelpaar und eine einzelne Modellantwort. Antworten von Sprachmodellen variieren zwischen Aufrufen, Modellen und Anbietern.
  \item \textbf{Reihenfolge.} Die Reihenfolge war mit 8:3 nicht ausbalanciert, sodass sich Reihenfolge- und Varianteneffekt nicht sauber trennen lassen.
  \item \textbf{Sprachfassungen.} Die deutsche Fassung von A ist eine maschinelle Übersetzung. Die beiden Fassungen von B unterscheiden sich in der Modalität: Nach der englischen Fassung erhöht die Kombination das Risiko (\enquote{significantly increases}), nach der deutschen \enquote{kann} sie es erheblich erhöhen. Welche Fassung die einzelnen Personen lasen, wurde nicht dokumentiert.
  \item \textbf{Datenqualität.} Bei K3 widerspricht ein Satz des Transkripts der Präferenz im Fragebogen. Ausgewertet wurde der Fragebogen, da der übrige Verlauf des Interviews ihn stützt. Bei K11 wurden vertauschte Bezeichnungen korrigiert, bei K8 eine missverstandene Frage ausgeschlossen. Die mündlichen Interviews wurden automatisch transkribiert und enthalten Erkennungsfehler.
  \item \textbf{Versuchsleitereffekt.} Entwicklung und Befragung lagen in derselben Hand, und die Teilnehmenden stammen aus dem persönlichen Umfeld. Sozial erwünschtes Antwortverhalten zugunsten von Variante~B ist daher nicht auszuschließen.
  \item \textbf{Umfang.} Bewertet wurde nur der Erklärungstext, nicht die Oberfläche mit Ampel und Detailansicht. Die Korrektheit der Risikoeinstufung wurde nicht gegen einen klinischen Referenzstandard wie Lexicomp geprüft.
\end{itemize}

\chapter{Konklusion}
\label{ch:konklusion}

\section{Zusammenfassung}

Diese Arbeit hat mit Medcheck einen Prototyp entwickelt und evaluiert, der amtliche Fachinformationen zu Arzneimittelwechselwirkungen mithilfe eines Sprachmodells für medizinische Laien aufbereitet. Ausgangspunkt war die Beobachtung, dass verlässliche Informationen zwar frei verfügbar, aber schwer verständlich sind, während frei antwortende Sprachmodelle zwar verständlich formulieren, Wechselwirkungen aber mit unzureichender Spezifität erkennen~\cite{MOHAMMED2026100733}. Medcheck verbindet beide Seiten, indem das Sprachmodell ausschließlich den amtlichen Text erklärt und selbst nicht als Wissensquelle dient.

\section{Beantwortung der Forschungsfragen}

\paragraph{FF1.} Die technische Realisierung gelingt mit einer Architektur, die Datenbeschaffung und Sprachverarbeitung strikt trennt. RxNorm normalisiert die Eingaben, openFDA liefert den Interaktionsabschnitt der Fachinformation, und ein Sprachmodell überführt ihn in eine Risikoeinstufung und eine Erklärung in einer von elf Sprachen. Drei Entwurfsentscheidungen erwiesen sich als tragend. Erstens kapselt das Interface \texttt{LlmProvider} die Anbieter, sodass Claude, Gemini und ein lokales Modell ohne Codeänderung austauschbar sind. Zweitens wird die Modellausgabe als nicht vertrauenswürdige Eingabe behandelt und bei Formatfehlern auf \texttt{UNKNOWN} zurückgeführt. Drittens bleibt durch die kontrollierte Degradation der amtliche Text auch dann verfügbar, wenn das Modell ausfällt. Der größte Aufwand lag dabei nicht im Sprachmodell, sondern in der robusten Anbindung der öffentlichen Datenquellen.

\paragraph{FF2.} Medizinische Laien bewerteten die von Medcheck erzeugte Erklärung im Mittel als verständlicher als den Text von Drugs.com (4{,}30 gegenüber 3{,}64), und acht von elf Personen fanden sie im direkten Vergleich verständlicher. Beim Vertrauen lagen beide Varianten nahezu gleichauf, bei der Nützlichkeit leicht zugunsten von Medcheck. Die Unterschiede sind bei elf Personen nicht signifikant und werden von einem möglichen Reihenfolgeeffekt überlagert. Inhaltlich zeigt sich ein Zielkonflikt: Die Kürze, die die Verständlichkeit erhöht, geht zulasten von Details wie konkreten Warnzeichen, die für Vertrauen und Risikoeinschätzung wichtig sind.

\section{Ausblick}

Aus den Ergebnissen ergeben sich folgende Ansatzpunkte für die Weiterentwicklung:

\begin{itemize}
  \item \textbf{Breitere Datengrundlage.} Neben dem Interaktionsabschnitt sollten die Warnhinweise und die Hinweise zur Patientenberatung der Fachinformation einbezogen werden. Die Zusammenfassung könnte dann konkrete Warnzeichen nennen, ohne das Grounding aufzugeben. Zudem sollte der Prompt verlangen, die Wirkrichtung eines Mechanismus ausdrücklich zu benennen.
  \item \textbf{Validierung der Einstufung.} Die Risikoeinstufung sollte wie bei Mohammed et al.~\cite{MOHAMMED2026100733} gegen einen Referenzstandard wie Lexicomp geprüft werden. Harmlose Kombinationen sollten dabei einbezogen werden, um die Spezifität der konservativen Regel zu bestimmen.
  \item \textbf{Zuordnung über Substanzkennungen.} Statt der namensbasierten Abfragekaskade könnte die Zuordnung zwischen RxNorm und openFDA über eindeutige Substanzkennungen wie den UNII-Code erfolgen.
  \item \textbf{Robustheit und Betrieb.} Wiederholungsversuche mit exponentiellem Backoff, automatisierte Tests der Degradation im \texttt{InteractionService}, Messungen der Antwortzeiten und eine persistente Datenbank mit Ablaufdatum für Cache-Einträge würden den Prototyp näher an einen produktiven Einsatz bringen.
  \item \textbf{Größere Studie.} Eine Folgestudie sollte die Reihenfolge ausbalancieren, ältere Menschen mit Polypharmazie einbeziehen, mehrere Arzneimittelpaare und Anbieter vergleichen und die vollständige Oberfläche einschließlich Ampel und Detailansicht bewerten.
\end{itemize}

Medcheck zeigt, dass ein Sprachmodell als Erklärschicht über amtlichen Daten Laien einen Mehrwert bieten kann, ohne selbst zur Wissensquelle zu werden. Die Ergebnisse legen nahe, dass die entscheidende Frage dabei weniger die Wahl des Modells ist als die Gestaltung der Information: welche Quelle zugrunde liegt, wie viel Detail gezeigt wird und wie deutlich das Risiko kommuniziert wird.


\printbibliography[heading=bibintoc, title={Literaturverzeichnis}]

\IfFileExists{kapitel/anhang-umfrage.tex}{%
  \appendix
  \input{kapitel/anhang-umfrage}%
}{}

\end{document}
